\chapter{Distance Sweep and Regime Dependence}

\section{Question}

After the Raw-RP objective was selected, the next question was whether its
benefit was tied to a single turbulence setting or whether it scaled across
distance. The \texttt{0405-distance-sweep-rawrp-f6p5mm} family fixes the same detector
objective and varies the propagation distance from 100\,m to 3000\,m.

\runfigure{0405-distance-sweep-rawrp-f6p5mm/21_distance_sweep_summary_6panel.png}
  {Distance sweep summary}
  {Six-panel summary of the 0405 distance sweep.}

\section{Three Regimes}

The weak, moderate, and strong labels are shorthand for the joint behavior of
$D/r_0$ and the Rytov variance $\sigma_R^2$. Small values indicate weak phase
perturbations, order-unity values indicate moderate scintillation, and large
values indicate a severely distorted detector input.

\begin{table}[H]
  \centering
  \begin{tabular}{rcccccc}
    \toprule
    \smalltablehead{$L$ (m)} & \smalltablehead{$D/r_0$} & \smalltablehead{$\sigma_R^2$} & \smalltablehead{Turb.\ bucket} & \smalltablehead{D2NN bucket} & \smalltablehead{Gain} & \smalltablehead{Regime} \\
    \midrule
    100 & 1.261 & 0.0146 & 382.10 & 728.41 & +90.63\% & weak \\
    500 & 3.312 & 0.2793 & 499.55 & 530.53 & +6.20\% & moderate \\
    1000 & 5.020 & 0.9955 & 166.01 & 180.91 & +8.98\% & moderate \\
    2000 & 7.609 & 3.5475 & 9.29 & 16.41 & +76.53\% & strong \\
    3000 & 9.704 & 7.4603 & 0.643 & 2.077 & +222.97\% & strong \\
    \bottomrule
  \end{tabular}
  \caption{Detector bucket power across the distance sweep. The relative gains
  at 2000--3000\,m are large because the turbulent baseline has already
  collapsed to a very small absolute value.}
\end{table}

\section{Weak and Moderate Turbulence}

At 100\,m, the baseline beam is still bright, but the D2NN can sharply increase
detector concentration. At 500--1000\,m, the gain is positive but modest,
showing that the system is no longer in a trivial beam-shaping regime and not
yet in a catastrophic collapse regime.

\section{Strong Turbulence}

At 2000--3000\,m the relative improvement becomes dramatic, but the absolute
numbers reveal why this must be interpreted carefully. A gain from 0.643 to
2.077 is detector-relevant, yet it is still far from the vacuum bucket level.
The D2NN is improving a badly collapsed detector signal, not restoring the
vacuum field.

\section{Cross-Distance Generalization}

\begin{table}[H]
  \centering
  \begin{tabular}{rcc}
    \toprule
    \smalltablehead{Distance} & \smalltablehead{3\,km model gain} & \smalltablehead{Local model gain} \\
    \midrule
    100\,m & -3.23\% & +90.63\% \\
    500\,m & -60.53\% & +6.20\% \\
    1000\,m & -70.49\% & +8.98\% \\
    2000\,m & -8.68\% & +76.53\% \\
    3000\,m & +222.97\% & +222.97\% \\
    \bottomrule
  \end{tabular}
  \caption{Cross-distance transfer of the 3\,km-trained model. Detector gains do
  not transfer reliably across distance regimes.}
\end{table}

\runfigure{0405-distance-sweep-rawrp-f6p5mm/22_cross_distance_3km_model/focal_L100m_3km_vs_local.png}
  {3 km model at 100 m}
  {Cross-distance comparison between the local and 3\,km models at 100\,m.}

\runfigure{0405-distance-sweep-rawrp-f6p5mm/22_cross_distance_3km_model/focal_L3000m_3km_vs_local.png}
  {3 km model at 3000 m}
  {The only fair match: the 3\,km model on the 3\,km data.}

\section{Interpretation}

The distance sweep shows that detector-side diffractive preprocessing is highly
regime-dependent. The optical strategy remains useful, but it does not define a
single universal front end. This motivates a comparison with the 4f FD2NN
architecture, which is taken up next.
